\begin{trapbox}

	\begin{description}[style=nextline, leftmargin=2.8em, labelsep=0.5em, itemsep=0.35em]
		\item[\textbf{\textcolor{CTrap}{易错点一（弧位误判）}}]
			学生常不加区分地使用$\angle ACB=\frac12\angle AOB$，忽略点$C$的位置。
		\item[\textbf{\textcolor{CTrap}{正确理解（分清优劣）：}}]
			根据$C$在优弧还是劣弧，选择$\angle ACB=\frac12\angle AOB$或$\angle ACB=\frac12(360^\circ-\angle AOB)$。
		\item[\textbf{\textcolor{CTrap}{错因分析（忽视分类）：}}]
			没有意识到圆周角定理包含两种情况，默认使用劣弧公式。
	\end{description}

	\medskip
	\begin{description}[style=nextline, leftmargin=2.8em, labelsep=0.5em, itemsep=0.35em]
		\item[\textbf{\textcolor{CTrap}{易错点二（前提遗忘）}}]
			使用等弧对等角时，忽略“在同圆或等圆中”这一前提。
		\item[\textbf{\textcolor{CTrap}{正确理解（同圆等圆）：}}]
			只有在同圆或等圆中，相等的弧所对的圆周角才相等。
		\item[\textbf{\textcolor{CTrap}{错因分析（条件缺失）：}}]
			误以为长度相等的弧就是等弧，没有考虑弧所在的圆必须相同。
	\end{description}

	\medskip
	\begin{description}[style=nextline, leftmargin=2.8em, labelsep=0.5em, itemsep=0.35em]
		\item[\textbf{\textcolor{CTrap}{易错点三（运算错误）}}]
			计算优弧圆周角时，忘记除以2，或计算$360^\circ-\angle AOB$后直接等于圆周角。
		\item[\textbf{\textcolor{CTrap}{正确理解（完整公式）：}}]
			必须严格代入$\angle ACB=\frac12(360^\circ-\angle AOB)$，先减后除。
		\item[\textbf{\textcolor{CTrap}{错因分析（顺序疏忽）：}}]
			混淆运算顺序，导致结果误差。
	\end{description}

\end{trapbox}
